\section{Error Assessment and Adaptive Snapshot Enrichment}
\label{sec:pod_enrichment}

The original POD basis provides a compact spatial representation
of the temperature fields in the training dataset. However, its
ability to represent solutions at parameter combinations outside
the training set must be assessed separately. This section
evaluates the original reduced-order model (ROM) over an
independent parameter sample, examines the effect of the number
of retained modes, and introduces a targeted snapshot-enrichment
procedure to improve the spatial basis.

The refinement procedure considers three approaches: increasing
the number of retained POD modes, identifying parameter regions
associated with relatively large ROM errors, and adding
full-order snapshots around selected high-error locations.
The resulting augmented basis is then compared with the original
basis.

\subsection{Initial ROM error assessment}

The reconstruction accuracy of a POD basis on its training
snapshots does not establish its predictive accuracy at new
parameter combinations. The ROM must also be evaluated at
parameter values that were not included in the construction
of the original snapshot matrix.

For this purpose, an independent set of 100 parameter
combinations is generated using Latin hypercube sampling (LHS)
over the prescribed parameter domain:
\[
1\leq\mu\leq100,
\qquad
0.01\leq Br\leq2.
\]
Unlike the uniformly spaced original training grid, the LHS
samples provide a separate distribution of parameter locations
throughout the domain.

For each sampled parameter pair $(\mu_k,Br_k)$, the full-order
model is solved to obtain a reference temperature field.
The corresponding ROM solution is then calculated using the
original POD basis. The resulting errors are recorded for
different numbers of retained modes.

This LHS dataset is used as an error-screening dataset. Its
purpose is to examine the performance of the original basis,
assess the effect of the retained basis dimension, and identify
parameter combinations that may require additional snapshots.
The 100 LHS cases are separate from the original training
snapshots and are not used to construct the initial POD basis.
They are also distinct from the subsequent dataset used for
the original-versus-augmented basis comparison.

\subsection{Error metrics}

The difference between the full-order and reduced-order
temperature fields is quantified using three error measures.
Let $\mathbf{T}_{\mathrm{FOM}}$ and
$\mathbf{T}_{\mathrm{ROM}}$ denote the corresponding
temperature vectors, each containing $N$ physical
temperature values.

The relative discrete $L_2$ error is defined as
\begin{equation}
    \boxed{
    \epsilon_{\mathrm{rel}}
    =
    \frac{
    \left\|
    \mathbf{T}_{\mathrm{FOM}}
    -
    \mathbf{T}_{\mathrm{ROM}}
    \right\|_2
    }{
    \left\|
    \mathbf{T}_{\mathrm{FOM}}
    \right\|_2
    }.
    }
    \label{eq:relative_l2_error}
\end{equation}
This measure quantifies the overall difference between the
two temperature fields relative to the magnitude of the
full-order solution.

The root-mean-square error (RMSE) is calculated as
\begin{equation}
    \boxed{
    RMSE
    =
    \sqrt{
    \frac{1}{N}
    \sum_{i=1}^{N}
    \left(
    T_{\mathrm{FOM},i}
    -
    T_{\mathrm{ROM},i}
    \right)^2
    }.
    }
    \label{eq:rom_rmse}
\end{equation}
The RMSE measures the average magnitude of the temperature
differences over all physical degrees of freedom.

The maximum absolute error is defined as
\begin{equation}
    \boxed{
    \epsilon_{\max}
    =
    \max_i
    \left|
    T_{\mathrm{FOM},i}
    -
    T_{\mathrm{ROM},i}
    \right|.
    }
    \label{eq:rom_max_error}
\end{equation}
Unlike the RMSE, this measure is determined by the single
spatial location with the largest absolute temperature
difference.

These measures provide complementary information. The
relative $L_2$ error describes the overall field discrepancy
relative to the full-order solution, the RMSE quantifies the
average absolute difference per degree of freedom, and the
maximum absolute error identifies the largest local
discrepancy.

All three metrics are calculated using the complete
reconstructed temperature field, rather than only the
temperature-deviation field. Thus, the comparison is performed
between $\mathbf{T}_{\mathrm{FOM}}$ and
$\mathbf{T}_{\mathrm{ROM}}$ after the reference temperature
has been added to the ROM reconstruction.

\subsection{Mode-count study}

The first refinement step examines the effect of increasing
the number of retained POD modes. The same original POD basis
is evaluated using the following mode counts:
\[
\boxed{
r\in\{1,2,3,4,5,6,8,10\}.
}
\]
For each selected value of $r$, the first $r$ modes of the
original basis are retained. The Galerkin reduced system is
then constructed and solved for each of the 100 LHS
parameter combinations.

The mean, median, 95th-percentile, and maximum relative
$L_2$ errors are reported in Table~\ref{tab:mode_count_errors}.
The mean RMSE is also included to provide information about
the absolute temperature discrepancy.

\begin{table}[htbp]
    \centering
    \caption{ROM error statistics for different numbers of
    retained POD modes over the 100 LHS screening cases.}
    \label{tab:mode_count_errors}
    \small
    \begin{tabular}{ccccccc}
        \toprule
        \textbf{Modes} &
        \textbf{Cases} &
        \textbf{Mean} &
        \textbf{Median} &
        \textbf{95th percentile} &
        \textbf{Maximum} &
        \textbf{Mean RMSE} \\
        \midrule
        1  & 100 & 0.395587 & 0.418291 & 0.614300 & 0.636144 & 0.813243 \\
        2  & 100 & 0.171552 & 0.175393 & 0.348644 & 0.410623 & 0.378429 \\
        3  & 100 & 0.027889 & 0.025590 & 0.063268 & 0.079345 & 0.073850 \\
        4  & 100 & 0.007663 & 0.006444 & 0.020409 & 0.044846 & 0.024056 \\
        5  & 100 & 0.002532 & 0.001145 & 0.013146 & 0.031450 & 0.010355 \\
        6  & 100 & 0.001377 & 0.000184 & 0.008551 & 0.024884 & 0.006698 \\
        8  & 100 & 0.000726 & $3.02\times10^{-6}$ & 0.005354 & 0.015513 & 0.003776 \\
        10 & 100 & 0.000468 & $2.96\times10^{-8}$ & 0.003338 & 0.010435 & 0.002440 \\
        \bottomrule
    \end{tabular}
\end{table}

The results show a consistent reduction in ROM error as
the number of retained modes increases. With one mode,
the mean relative $L_2$ error is approximately 39.56\%,
and the maximum error reaches 63.61\%. Increasing the
basis dimension to three modes reduces the mean error
to 2.79\%, while four modes reduce it further to 0.77\%.

The improvement continues with additional modes. For
ten retained modes, the mean relative $L_2$ error
decreases to approximately 0.0468\%, and the median
reaches $2.96\times10^{-8}$. However, the 95th-percentile
error remains approximately 0.334\%, and the maximum
error is approximately 1.04\%. This difference between
the median and upper-tail statistics indicates that
the ROM errors are not uniformly distributed across
the test cases. Although the typical error becomes very
small at higher basis dimensions, some parameter
combinations still exhibit comparatively larger errors.

The maximum error does not necessarily occur at the same
parameter combination for every mode count. The locations
associated with the maximum errors are reported in
Table~\ref{tab:mode_count_error_locations}.

\begin{table}[htbp]
    \centering
    \caption{Parameter combinations associated with the
    maximum relative $L_2$ error at each retained mode count.}
    \label{tab:mode_count_error_locations}
    \small
    \begin{tabular}{cccc}
        \toprule
        \textbf{Modes} &
        \textbf{Case ID} &
        $\boldsymbol{\mu}$ &
        $\boldsymbol{Br}$ \\
        \midrule
        1  & \texttt{lhs\_test\_0054} & 92.777390 & 1.948657 \\
        2  & \texttt{lhs\_test\_0016} & 98.197590 & 1.610037 \\
        3  & \texttt{lhs\_test\_0070} & 5.720724 & 1.320725 \\
        4  & \texttt{lhs\_test\_0075} & 4.066653 & 1.769485 \\
        5  & \texttt{lhs\_test\_0075} & 4.066653 & 1.769485 \\
        6  & \texttt{lhs\_test\_0075} & 4.066653 & 1.769485 \\
        8  & \texttt{lhs\_test\_0075} & 4.066653 & 1.769485 \\
        10 & \texttt{lhs\_test\_0099} & 3.035937 & 0.320114 \\
        \bottomrule
    \end{tabular}
\end{table}

The maximum-error location changes with the retained
basis dimension. For three modes, the largest error
occurs at $(\mu,Br)=(5.720724,1.320725)$. For four,
five, six, and eight modes, the same parameter
combination, $(4.066653,1.769485)$, produces the
maximum error. At ten modes, the maximum shifts to
$(3.035937,0.320114)$. This variation indicates that
increasing the number of modes does not reduce the
error uniformly at every parameter location.

A plot of the mean, median, 95th-percentile, and maximum
relative $L_2$ errors against the retained mode count
provides a visual comparison of the error distributions.
A logarithmic vertical axis is suitable because the
error values span several orders of magnitude.

\subsection{Identification of high-error parameter regions}

The LHS screening results are also used to identify
parameter combinations for which the original ROM
produces relatively large errors. Each evaluation case
is associated with its parameter coordinates,
\[
(\mu_k,Br_k),
\]
and its corresponding relative $L_2$ error,
$\epsilon_{\mathrm{rel},k}$.

The error values are examined over the two-dimensional
parameter space. A scatter plot of the LHS samples,
coloured according to their relative $L_2$ errors,
provides a spatial representation of the ROM accuracy
throughout the sampled domain.

This visualization helps identify regions where the
original basis produces larger errors. These regions
provide candidate locations for additional training
snapshots. The parameter coordinates are retained
alongside the error values so that the high-error
cases can be ranked and grouped during the subsequent
hotspot-selection procedure.

The errors are not necessarily distributed uniformly
across the parameter domain. Consequently, simply
increasing the number of uniformly spaced training
points may not efficiently target the locations where
the original reduced basis has the greatest difficulty.
The enrichment procedure instead uses the case-wise
ROM errors to guide the selection of additional
snapshot locations.

\subsection{Selection of hotspot candidates}

The hotspot-selection procedure considers the results
from four selected mode counts:
\[
\boxed{
r\in\{5,6,8,10\}.
}
\]
For each of these mode counts, the cases with the
largest ROM errors are considered. The mode-count
study provides the 20 highest-error records for each
selected mode count. From these records, the six
highest-ranked cases are selected as candidates.

Candidates are ranked separately within each mode
count because the magnitude and distribution of the
errors change with the number of retained modes.
Therefore, raw error values are not compared directly
across different mode counts. Instead, each candidate
receives an order score according to its position in
the error ranking for the corresponding mode count.
A higher rank score represents a more severe position
in that ranking.

The candidate records from the different mode counts
are then grouped according to their parameter
combinations $(\mu,Br)$. This allows a parameter
combination that appears among the high-error cases
for multiple mode counts to be identified and given
additional support.

Two quantities are calculated for each grouped
parameter combination.

First, the mean rank score, $S_{\mathrm{rank}}$,
represents its average ranking across the mode counts
in which it appears. A higher score indicates that
the parameter combination is consistently associated
with a relatively large ROM error.

Second, the support fraction, $S_{\mathrm{support}}$,
measures how many of the selected mode counts identify
the same parameter combination as a high-error
candidate. It is defined as
\begin{equation}
    \boxed{
    S_{\mathrm{support}}
    =
    \frac{n_{\mathrm{mode}}}{N_{\mathrm{mode}}},
    }
    \label{eq:hotspot_support}
\end{equation}
where $n_{\mathrm{mode}}$ is the number of selected
mode counts in which the parameter combination appears,
and $N_{\mathrm{mode}}=4$ is the total number of
mode counts considered.

The two quantities are combined to obtain the base
candidate score:
\begin{equation}
    \boxed{
    S_{\mathrm{base}}
    =
    0.70S_{\mathrm{rank}}
    +
    0.30S_{\mathrm{support}}.
    }
    \label{eq:hotspot_base_score}
\end{equation}
This weighting assigns 70\% of the score to the
candidate's mean ranking and 30\% to its support
across different mode counts. The resulting base
scores provide an initial ranking of the hotspot
candidates.

However, selecting candidates solely according to
their base scores could produce several locations
close to one another in parameter space. A further
selection criterion is therefore introduced to
promote spatial diversity.

\subsection{Spatially diverse hotspot selection}

Selecting the highest-scoring candidates independently
may lead to several hotspots concentrated in a
relatively small region of the parameter domain.
Such a selection would provide limited coverage of
other potentially problematic regions.

To promote spatial diversity, the hotspot-selection
algorithm uses a normalized Euclidean distance in
the two-dimensional parameter space. The normalized
parameter coordinates are defined as
\begin{equation}
    \tilde{\mu}
    =
    \frac{\mu-\mu_{\min}}
    {\mu_{\max}-\mu_{\min}},
    \qquad
    \widetilde{Br}
    =
    \frac{Br-Br_{\min}}
    {Br_{\max}-Br_{\min}}.
    \label{eq:normalized_parameter_coordinates}
\end{equation}

The normalized distance between two parameter
combinations $p$ and $q$ is
\begin{equation}
    d(p,q)
    =
    \sqrt{
    (\tilde{\mu}_p-\tilde{\mu}_q)^2
    +
    (\widetilde{Br}_p-\widetilde{Br}_q)^2
    }.
    \label{eq:normalized_parameter_distance}
\end{equation}

At each selection step, the minimum distance from
an unselected candidate $p$ to the already selected
hotspots is calculated:
\begin{equation}
    \boxed{
    d_{\min}(p)
    =
    \min_{q\in\mathcal{H}}d(p,q),
    }
    \label{eq:minimum_hotspot_distance}
\end{equation}
where $\mathcal{H}$ denotes the set of selected
hotspot centres.

The first hotspot is selected according to the
highest base score. For subsequent selections, the
candidate score combines its base score with its
minimum distance from the existing hotspot centres:
\begin{equation}
    \boxed{
    S_{\mathrm{selection}}
    =
    0.55S_{\mathrm{base}}
    +
    0.45d_{\min}.
    }
    \label{eq:hotspot_selection_score}
\end{equation}
The candidate with the highest updated selection
score is added to the hotspot set. Its inclusion
changes the minimum-distance values of the remaining
candidates, and the procedure is repeated until the
maximum permitted number of hotspot centres is
reached or the available candidates are exhausted.

The algorithm permits up to six hotspot centres.
However, the actual selection procedure produced
three centres for the present dataset. Their
coordinates are reported in
Table~\ref{tab:selected_hotspots}.

\begin{table}[htbp]
    \centering
    \caption{Selected hotspot centres in the parameter domain.}
    \label{tab:selected_hotspots}
    \begin{tabular}{ccc}
        \toprule
        \textbf{Hotspot} &
        $\boldsymbol{\mu}$ &
        $\boldsymbol{Br}$ \\
        \midrule
        1 & 4.066653 & 1.769485 \\
        2 & 3.035937 & 0.320114 \\
        3 & 5.720724 & 1.320725 \\
        \bottomrule
    \end{tabular}
\end{table}

The selected locations correspond to parameter
combinations that appear among the maximum-error
cases for several of the tested mode counts. In
particular, the first hotspot is the maximum-error
location for four, five, six, and eight modes. The
third hotspot is the maximum-error location for
three modes, while the second is the maximum-error
location for ten modes. This agreement provides
additional motivation for using these locations
as centres for local enrichment.

A parameter-space scatter plot can be used to show
the LHS screening samples and the selected hotspot
centres. This provides a visual representation of
the selected enrichment locations relative to the
parameter domain.

\subsection{Local parameter-space refinement}

Additional parameter combinations are generated
around each selected hotspot to enrich the original
snapshot dataset in regions associated with
relatively large ROM errors.

The original uniformly spaced training grid has
parameter spacings
\begin{equation}
    \Delta\mu
    =
    \frac{100-1}{10-1}
    =11,
\end{equation}
and
\begin{equation}
    \Delta Br
    =
    \frac{2-0.01}{10-1}
    \approx0.221111.
\end{equation}

For each hotspot centre $(\mu_h,Br_h)$, local
parameter combinations are generated using the
offset coefficients
\[
\delta_\mu,\delta_{Br}
\in\{-0.5,0,0.5\}.
\]
The corresponding parameter values are
\begin{equation}
    \mu
    =
    \mu_h+\delta_\mu\Delta\mu,
    \qquad
    Br
    =
    Br_h+\delta_{Br}\Delta Br.
    \label{eq:local_parameter_refinement}
\end{equation}

The Cartesian product of these offsets defines a
local $3\times3$ candidate grid around each hotspot.
The centre itself is excluded because it is already
represented by the hotspot location. The remaining
candidates are filtered according to the following
conditions:
\begin{itemize}
    \item Parameter combinations outside the original
    parameter domain are rejected.
    \item Duplicate parameter combinations are removed.
    \item Combinations already present in the original
    training grid are excluded.
\end{itemize}

After applying these restrictions to the three
selected hotspot centres, 15 additional parameter
combinations remain.

Each new parameter combination is passed to the
full-order solver. Its converged temperature field
is then processed using the same snapshot-preparation
procedure as the original training data. Ghost cells
are removed, the reference temperature is subtracted,
and the resulting deviation field is flattened and
stored as an additional snapshot column.

The local refinement therefore increases the
sampling density around the selected high-error
regions without uniformly increasing the number
of training points throughout the entire parameter
domain.

\subsection{Construction of the augmented POD basis}

The augmented snapshot matrix is constructed by
combining the original snapshots with the additional
snapshots generated through local refinement:
\begin{equation}
    \boxed{
    \mathbf{X}_{\mathrm{aug}}
    =
    \begin{bmatrix}
    \mathbf{X}_{\mathrm{original}}
    &
    \mathbf{X}_{\mathrm{additional}}
    \end{bmatrix}.
    }
    \label{eq:augmented_snapshot_matrix}
\end{equation}

The dimensions of the two snapshot datasets are
\[
\mathbf{X}_{\mathrm{original}}
\in\mathbb{R}^{1000\times100},
\]
and
\[
\mathbf{X}_{\mathrm{additional}}
\in\mathbb{R}^{1000\times15}.
\]
Consequently, the augmented snapshot matrix has
dimensions
\begin{equation}
    \boxed{
    \mathbf{X}_{\mathrm{aug}}
    \in\mathbb{R}^{1000\times115}.
    }
    \label{eq:augmented_snapshot_dimensions}
\end{equation}

A new SVD is performed on the augmented snapshot
matrix:
\begin{equation}
    \mathbf{X}_{\mathrm{aug}}
    =
    \mathbf{U}_{\mathrm{aug}}
    \mathbf{\Sigma}_{\mathrm{aug}}
    \mathbf{W}_{\mathrm{aug}}^T.
    \label{eq:augmented_svd}
\end{equation}
The augmented POD basis is obtained by retaining
the first $r$ columns of
$\mathbf{U}_{\mathrm{aug}}$:
\begin{equation}
    \boxed{
    \mathbf{V}_{r,\mathrm{aug}}
    =
    \mathbf{U}_{\mathrm{aug}}(:,1:r).
    }
    \label{eq:augmented_pod_basis}
\end{equation}

Adding the additional snapshots changes the snapshot
matrix and therefore changes its singular values
and spatial modes. The cumulative energy captured
by the first three modes of the original and
augmented bases is compared in
Table~\ref{tab:original_augmented_energy}.

\begin{table}[htbp]
    \centering
    \caption{Comparison of three-mode cumulative POD
    energy for the original and augmented snapshot datasets.}
    \label{tab:original_augmented_energy}
    \begin{tabular}{ccc}
        \toprule
        \textbf{Basis} &
        \textbf{Number of snapshots} &
        \textbf{Cumulative energy in first 3 modes (\%)} \\
        \midrule
        Original & 100 & 99.9833 \\
        Augmented & 115 & 99.9699 \\
        \bottomrule
    \end{tabular}
\end{table}

The first three modes of the original basis capture
approximately 99.9833\% of its snapshot energy,
whereas the first three modes of the augmented basis
capture approximately 99.9699\%. The slight decrease
in cumulative energy reflects the redistribution of
the snapshot energy after the additional temperature
fields are included.

This difference does not, by itself, establish
whether the original or augmented basis gives more
accurate predictions at unseen parameter
combinations. The POD energy is calculated relative
to the total energy of its own snapshot matrix.
Adding snapshots changes that reference total and
the resulting singular-value distribution. A lower
cumulative energy at a fixed number of modes
therefore does not necessarily correspond to a
larger prediction error.

The effect of snapshot enrichment must instead be
assessed by evaluating both bases on the same
independent parameter set, using the same retained
mode counts and error metrics. This comparison is
presented in the subsequent results section.

\subsection{Code implementation}
\label{subsec:refinement_code}

The refinement and augmented-basis construction
workflow is implemented through six Python scripts
in the \texttt{F05\_refinement} directory.

\begin{itemize}

    \item \texttt{f00\_generate\_large\_training\_cases.py}:
    Generates the 100-case training parameter grid
    over the prescribed parameter domain. The
    function \texttt{build\_parameter\_grid} creates
    the parameter combinations and saves the grid
    as a CSV file. The script then processes each
    case using \texttt{run\_fom\_case}. Converged
    temperature solutions are converted into
    snapshot columns, and their associated
    simulation information is stored in
    \texttt{snapshot\_case\_records}. The snapshot
    columns and case records are saved for
    subsequent POD construction.

    \item \texttt{f01\_build\_pod\_basis.py}:
    Loads the training snapshot matrix and its
    associated case records, performs the SVD,
    and saves the resulting POD data. This script
    constructs the original basis used in the
    subsequent mode-count study.

    \item \texttt{f02\_generate\_large\_test\_cases.py}:
    Generates the LHS evaluation dataset. Its
    function \texttt{build\_latin\_hypercube\_cases}
    creates the parameter combinations within the
    prescribed domain. The function
    \texttt{run\_lhs\_test\_case\_generation}
    then runs the full-order solver for these
    combinations and saves the resulting case
    records.

    \item \texttt{f03\_mode\_count\_study.py}:
    Evaluates the original POD basis using the
    selected retained-mode counts. The function
    \texttt{run\_rom\_for\_mode\_count} performs
    the Galerkin projection and reduced-order
    solution for a specified case and mode count.
    The main function, \texttt{run\_mode\_count\_study},
    loads the POD modes and LHS evaluation cases,
    constructs the required full-order operators,
    and evaluates the ROM for each selected mode
    count. It saves the case-wise error records
    and generates error visualizations. The
    function \texttt{plot\_error\_statistics}
    summarizes the error metrics as a function
    of the retained basis dimension.

    \item \texttt{f04\_generate\_augmented\_training\_set.py}:
    Selects the hotspot candidates and generates
    the additional training parameter combinations.
    The function \texttt{build\_hotspot\_candidates}
    loads the high-error cases from the mode-count
    study, considers the 20 worst-error records
    for each selected mode count, and selects and
    ranks the top candidates. It then groups
    repeated parameter combinations and calculates
    the mean rank score, support fraction, and
    base score for each group.

    The function \texttt{select\_hotspots} applies
    the spatial-diversity criterion to select the
    hotspot centres. The function
    \texttt{build\_augmented\_parameter\_grid}
    generates local parameter combinations around
    these centres and applies the parameter-domain,
    duplicate, and original-grid restrictions.
    Finally,
    \texttt{run\_augmented\_training\_set\_generation}
    solves the additional cases using the
    full-order model, adds their converged
    temperature fields to the augmented snapshot
    dataset, and saves the corresponding records.

    \item \texttt{f05\_build\_augmented\_pod\_basis.py}:
    Constructs the augmented POD basis. The
    function \texttt{combine\_snapshot\_datasets}
    loads the original and additional snapshot
    matrices and concatenates them column-wise.
    The function \texttt{build\_augmented\_pod\_basis}
    then performs the SVD on the combined matrix
    and saves the augmented POD results.

\end{itemize}

The scripts separate the error-screening process
from the generation of additional training data
and the construction of the augmented basis. The
original basis is first evaluated over the LHS
parameter set. Its high-error cases guide the
selection of local enrichment regions. Additional
full-order snapshots are then generated around
these regions and combined with the original
training data to obtain the augmented POD basis.